\documentclass[12pt]{article}
\usepackage{fullpage}
\usepackage[T1]{fontenc}
\usepackage[utf8]{inputenc}
\usepackage{lmodern}
\usepackage{microtype}
\usepackage{amsmath,amssymb,amsthm}
\usepackage{hyperref}

\newtheorem{theorem}{Theorem}
\newtheorem{proposition}{Proposition}
\newtheorem{lemma}{Lemma}
\newtheorem{corollary}{Corollary}
\theoremstyle{definition}
\newtheorem{definition}{Definition}
\newtheorem{remark}{Remark}

\newcommand{\Z}{\mathbb{Z}}
\newcommand{\R}{\mathbb{R}}
\newcommand{\N}{\mathbb{N}}
\newcommand{\E}{\mathbb{E}}
\newcommand{\Prob}{\mathbb{P}}

\begin{document}

\title{Generalized Zeckendorf Decompositions:\\ Existence/Uniqueness Phase Diagram, Gap Distribution, and Largest Gap}
\author{}
\date{}
\maketitle

\section{Setup, conventions, and the legality automaton}

Throughout, $L\ge1$ is fixed and $c_1,\dots,c_L$ are non-negative integers with
$c_1\ge1$, $c_L\ge1$. The \emph{positive linear recurrence sequence} (PLRS)
$\{H_n\}_{n\ge1}$ is defined by
\begin{equation}\label{eq:rec}
H_{n+1}=c_1H_n+c_2H_{n-1}+\cdots+c_LH_{n-L+1}\qquad(n\ge L),
\end{equation}
with the standard Miller--Wang initial conditions
\begin{equation}\label{eq:init}
H_1=1,\qquad H_n=c_1H_{n-1}+c_2H_{n-2}+\cdots+c_{n-1}H_1+1\quad(2\le n\le L).
\end{equation}
We write $a^+=\max(a,0)$ and use the convention that an empty sum is $0$.

\subsection{Legal decompositions}
A \emph{decomposition} of $N\in\Z_{\ge1}$ is an expression
$N=\sum_{i\ge1}a_iH_i$ with $a_i\in\Z_{\ge0}$ and only finitely many $a_i\neq0$.
Its \emph{digit string} is the finite word $a_m a_{m-1}\cdots a_1$, where
$m$ is the largest index with $a_m\neq0$ (read most-significant-first).

\begin{definition}[Legal decomposition; Miller--Wang]\label{def:legal}
A digit string $a_m\cdots a_1$ ($a_m\ge1$) is \emph{legal} if it does not
contain, starting at any position, the lexicographic pattern
$(c_1,c_2,\dots,c_{s-1},x)$ with $x\ge c_s$ as a leading block. Precisely,
$a_m\cdots a_1$ is legal iff it admits a factorization into consecutive
\emph{blocks}, read left to right, where each block is one of
\begin{equation}\label{eq:blocks}
(c_1,c_2,\dots,c_{j-1},b),\qquad 1\le j\le L,\ \ 0\le b<c_j ,
\end{equation}
and possibly a final \emph{incomplete} block $(c_1,\dots,c_t)$ with $t<L$
(a strict prefix of $(c_1,\dots,c_L)$). A decomposition is legal iff its digit
string is legal.
\end{definition}

For $L=2,c_1=c_2=1$ Definition~\ref{def:legal} is exactly Zeckendorf's rule
(no two consecutive Fibonacci numbers).

\subsection{The legality automaton}\label{ss:automaton}
We encode legality by a deterministic finite automaton $\mathcal A$ reading the
digit string most-significant-first. The states are $\{0,1,\dots,L-1\}$; state
$j$ means ``the current block has matched $c_1,\dots,c_j$ exactly so far''
(state $0$ marks the start of a fresh block). From state $j$ the next digit $a$
acts as follows, where we abbreviate $c_{j+1}$ for the relevant coefficient:
\begin{itemize}
\item[(C)] \emph{Continue.} If $a=c_{j+1}$ and $j+1<L$: move to state $j+1$.
\item[(T)] \emph{Terminate.} If $0\le a<c_{j+1}$: the block ends with deficient
last digit $a$; the next block begins, so move to state $0$.
\item[(F)] If $j+1=L$ and $a=c_L$, or if $a>c_{j+1}$: \emph{forbidden}.
\end{itemize}
A word is legal iff $\mathcal A$ reads it without entering (F), starting and
ending in admissible states; an incomplete final block corresponds to ending in
a state $t<L$.

\begin{lemma}[Transfer matrix]\label{lem:transfer}
Let $M=(M_{j'j})_{0\le j,j'\le L-1}$ be the integer matrix counting transitions
of $\mathcal A$:
\begin{equation}\label{eq:M}
M_{j+1,\,j}=1\ \ (0\le j\le L-2),\qquad M_{0,\,j}=c_{j+1}\ \ (0\le j\le L-1),
\end{equation}
all other entries $0$. Then for every $n\ge1$,
\begin{equation}\label{eq:countH}
\#\{\text{legal strings of length }n\}=H_{n+1},
\end{equation}
and the number of legal strings whose leading digit lies at index $n$ (i.e.\ of
$N\in[H_n,H_{n+1})$) equals $H_{n+1}-H_n=\sum_{j=0}^{L-1}M_{0,j}\,(\,\cdot\,)$,
which is the number of integers in $[H_n,H_{n+1})$.
\end{lemma}

\begin{proof}
Entry \eqref{eq:M} encodes (C) (the single continuation digit $a=c_{j+1}$
sending $j\mapsto j+1$) and (T) (the $c_{j+1}$ terminating digits
$a\in\{0,\dots,c_{j+1}-1\}$ sending $j\mapsto 0$); (F) contributes nothing.
Hence $(M^n)_{\bullet,0}$ counts, by destination state, the legal words of
length $n$ that begin a fresh block; summing over destinations and starting from
state $0$ gives the total count of legal length-$n$ words,
$\mathbf 1^{\!\top}M^{n}e_0$, where $e_0$ is the first standard basis vector.

We prove $g_n:=\mathbf 1^{\!\top}M^{n}e_0=H_{n+1}$ by induction. The
characteristic polynomial of $M$ is, expanding along the first row of
$xI-M$,
\begin{equation}\label{eq:char}
\chi(x)=x^{L}-c_1x^{L-1}-c_2x^{L-2}-\cdots-c_L ,
\end{equation}
so by the Cayley--Hamilton theorem each entry of $M^n$, hence $g_n$, satisfies
the recurrence \eqref{eq:rec}. A direct check using \eqref{eq:M} gives
$g_n=H_{n+1}$ for $1\le n\le L$ using \eqref{eq:init}, and then \eqref{eq:rec}
propagates the identity for all $n$. Finally, the number of $N\in[H_n,H_{n+1})$
is $H_{n+1}-H_n=g_n-g_{n-1}$, matching the count of legal strings with leading
digit exactly at index $n$.
\end{proof}

\begin{proposition}[Existence and uniqueness in the PLRS regime]\label{prop:eu}
For a PLRS with $c_1\ge1$, $c_L\ge1$, every positive integer has a unique legal
decomposition, obtained by the greedy algorithm.
\end{proposition}

\begin{proof}
Greedy algorithm: given $N$, let $m$ be maximal with $H_m\le N$, set
$a_m=\lfloor N/H_m\rfloor$ and recurse on $N-a_mH_m$ with indices $<m$.

\emph{Legality of the greedy output.} We show the produced string is read by
$\mathcal A$ without hitting (F). After choosing $a_m,\dots,a_{m-j+1}$ equal to
$c_1,\dots,c_j$ (state $j$), the remaining value satisfies
$N-\sum_{i=1}^{j}c_iH_{m-i+1}<H_{m-j+1}$ when $j<L$ (else by \eqref{eq:rec} the
prefix $(c_1,\dots,c_{j})$ would already account for $H_{m+1}>N$); hence the next
greedy digit is $<c_{j+1}$, i.e.\ a (T) or a further (C) move, never $a>c_{j+1}$
and never $a=c_L$ at $j+1=L$. Thus (F) never occurs.

\emph{Uniqueness.} Legal strings of length $n$ are in bijection (by
$\mathcal A$) with paths counted by Lemma~\ref{lem:transfer}, and
\eqref{eq:countH} shows their number equals $\#[H_n,H_{n+1})$. Distinct legal
strings represent distinct integers because for legal strings
$a_m\cdots a_1$ the value $\sum a_iH_i$ lies in $[H_m,H_{m+1})$ and the map is a
counting bijection onto $[H_m,H_{m+1})$ by the equality of cardinalities. Hence
each $N$ has exactly one legal decomposition.
\end{proof}

\section{Part (a): Relaxed legality regimes}\label{sec:a}

We now describe how existence/uniqueness, the mean number of summands, and the
Gaussian fluctuations are modified when the legality rule is relaxed. We make
precise three regimes by altering the per-state \emph{forbidden} set (F) of
Section~\ref{ss:automaton}.

\subsection{The reference (PLRS) regime}
For comparison we first record the reference asymptotics. Let $\lambda>1$ be the
unique dominant (Perron) root of \eqref{eq:char}; positivity of the coefficients
and $c_1,c_L\ge1$ make $M$ primitive
($\gcd\{i: c_i>0\}$ together with $c_1>0$ gives aperiodicity), so by the
Perron--Frobenius theorem $\lambda$ is real, simple, and strictly larger in
modulus than every other eigenvalue.

\begin{theorem}[Reference Gaussian behavior]\label{thm:gauss}
Let $K(N)$ be the number of summands (number of nonzero digits) in the unique
legal decomposition of $N$. If $N$ is uniform in $[H_n,H_{n+1})$, then there are
constants $\mu>0$, $\sigma^2>0$ depending only on $(c_1,\dots,c_L)$ such that
\begin{equation}\label{eq:mean}
\E[K]=\mu n+O(1),\qquad \mathrm{Var}(K)=\sigma^2 n+O(1),
\end{equation}
and $(K-\mu n)/\sqrt{\sigma^2 n}$ converges in distribution to the standard
normal $\mathcal N(0,1)$. Explicitly,
\begin{equation}\label{eq:density}
\mu=\beta:=1-\frac1S\sum_{j=0}^{L-1} w_j\,\lambda^{-(j+1)},\qquad
S=L-\sum_{i=1}^{L-1}(L-i)\,c_i\,\lambda^{-i},
\end{equation}
where $w_j=1$ if $c_{j+1}\ge1$ and $w_j=r_{j+1}$ if $c_{j+1}=0$, and
$r_j=\lambda^{j}-\sum_{i=1}^{j}c_i\lambda^{j-i}$ (the left Perron eigenvector of
$M$, see Lemma~\ref{lem:eig}).
\end{theorem}

\begin{proof}[Proof sketch with full ingredients]
By Lemma~\ref{lem:transfer} legal strings of length $n$ are exactly the
length-$n$ paths of $\mathcal A$, a primitive automaton. Mark each edge by a
formal variable $y$ when it emits a \emph{nonzero} digit (a summand) and by $1$
otherwise. The bivariate transfer matrix $M(y)$ has $M(y)_{j+1,j}=1$ if
$c_{j+1}=0$ and $=y$ if $c_{j+1}\ge1$; and
$M(y)_{0,j}=1+(c_{j+1}-1)y$ (one zero-terminating digit $a=0$ and
$c_{j+1}-1$ nonzero terminating digits). The generating function of $K$ over
length-$n$ words is $\mathbf 1^{\!\top}M(y)^n e_0$. By Perron--Frobenius applied
to $M(y)$ for $y$ near $1$, the dominant eigenvalue $\lambda(y)$ is analytic and
simple, so $K$ is, up to negligible spectral-gap corrections, a sum of
asymptotically independent contributions across blocks; the Quasi-Power Theorem
of Hwang (which requires exactly: analyticity of $\lambda(y)$ near $y=1$, a
spectral gap, and $\frac{d^2}{dy^2}\log\lambda(y)|_{y=1}\neq0$, all of which hold
here by primitivity and the fact that block lengths are bounded by $L$) yields
\eqref{eq:mean} with
$\mu=\lambda'(1)/\lambda(1)$, $\sigma^2=\frac{d^2}{dy^2}\log\lambda(y)|_{y=1}$,
and the Gaussian limit. The closed form \eqref{eq:density} for $\mu$ is the
stationary probability of a nonzero emission, computed in
Lemma~\ref{lem:base}.
\end{proof}

\subsection{Multiple decompositions: the over-determined regime}
Relax (F) by permitting digits up to $c_{j+1}$ even at $j+1=L$, i.e.\ allow the
full block $(c_1,\dots,c_L)$ to appear without being forced to carry. Equivalently
the legal alphabet is widened so that the constraint becomes
$a_i\le c_1$ with the carrying rule \eqref{eq:rec} not enforced as an
equality-avoidance. Then some integers acquire several legal decompositions
because the single relation
\begin{equation}\label{eq:relation}
H_{n+1}=c_1H_n+c_2H_{n-1}+\cdots+c_LH_{n-L+1}
\end{equation}
can be applied or not.

\begin{proposition}[Over-determined regime]\label{prop:multi}
In the relaxed automaton $\mathcal A^{+}$ in which transition (F) at
$j+1=L,\ a=c_L$ is replaced by a legal move to state $0$, every positive integer
has \emph{at least one} legal decomposition, and the number of legal
decompositions of $N$ uniform in $[H_n,H_{n+1})$ has mean growing geometrically:
there is $\rho>1$ (the spectral radius of the augmented transfer matrix
$M^{+}=M+e_0e_{L-1}^{\!\top}$ counting the extra edge) with
\[
\E[\#\text{decompositions of }N]\;\asymp\;(\rho/\lambda)^{\,n}.
\]
The mean number of summands per decomposition still satisfies
$\E[K]=\mu^{+}n+O(1)$ and obeys a central limit theorem with variance linear in
$n$; the constants $(\mu^{+},(\sigma^{+})^{2})$ are
$\mu^{+}=\frac{d}{dy}\log\rho(y)|_{1}$,
$(\sigma^{+})^{2}=\frac{d^2}{dy^2}\log\rho(y)|_{1}$, computed from the bivariate
version $M^{+}(y)$ exactly as in Theorem~\ref{thm:gauss}.
\end{proposition}

\begin{proof}
Existence: $\mathcal A^{+}\supseteq\mathcal A$, so the (unique) greedy string is
still legal, giving at least one decomposition. The number of legal length-$n$
words equals $\mathbf 1^{\!\top}(M^{+})^n e_0$; since $M^{+}\ge M$ entrywise with
a strictly larger primitive part (the extra edge $e_0e_{L-1}^{\!\top}$ adds a
cycle through state $L-1$), Perron--Frobenius gives spectral radius
$\rho>\lambda$, whence the number of legal length-$n$ words is $\Theta(\rho^n)$
while the number of integers represented is $\Theta(\lambda^n)$ (still
$[H_n,H_{n+1})$); the mean multiplicity is the ratio, $\Theta((\rho/\lambda)^n)$.
$M^{+}$ is primitive (it contains the primitive $M$ as a subgraph plus an edge),
so the bivariate $M^{+}(y)$ again has an analytic simple dominant eigenvalue
near $y=1$, and the Quasi-Power Theorem yields the linear-mean/linear-variance
Gaussian law for $K$. (Uniqueness fails precisely because the extra edge encodes
the optional use of relation \eqref{eq:relation}.)
\end{proof}

\subsection{No decomposition: the under-determined regime}
Now relax in the opposite direction: \emph{shrink} the legal alphabet so that
some terminating digits are forbidden, e.g.\ require the last digit of every
block to be $0$ (only the patterns $(c_1,\dots,c_{j-1},0)$ are legal blocks).
Then the represented set is a proper subset of $\Z_{\ge1}$.

\begin{proposition}[Under-determined regime]\label{prop:none}
Let $\mathcal A^{-}$ be the automaton obtained from $\mathcal A$ by deleting the
nonzero terminating digits (keep only $a=0$ in (T) and the single continuation
$a=c_{j+1}$ in (C)). Then the set $\mathcal R_n$ of integers in $[H_n,H_{n+1})$
admitting an $\mathcal A^{-}$-legal decomposition has size
$\#\mathcal R_n=\Theta(\rho_-^{\,n})$ with $1<\rho_-<\lambda$ (the spectral
radius of the reduced transfer matrix $M^{-}$), so
\[
\frac{\#\mathcal R_n}{\#[H_n,H_{n+1})}=\Theta\!\big((\rho_-/\lambda)^n\big)\to0 ,
\]
i.e.\ asymptotically almost no integer is representable. On the representable set
$\mathcal R_n$ (conditioned on existence), decompositions are unique, the mean
number of summands is $\mu^{-}n+O(1)$, and a central limit theorem with variance
$\Theta(n)$ holds, with constants computed from $M^{-}(y)$ as before.
\end{proposition}

\begin{proof}
Deleting edges from a primitive nonnegative matrix strictly decreases the
spectral radius (if the deleted edges lie on some cycle, which they do here:
nonzero terminating digits create cycles returning to state $0$), so
$1<\rho_-<\lambda$. The count $\#\mathcal R_n=\mathbf 1^{\!\top}(M^{-})^n e_0$
is $\Theta(\rho_-^n)$ by Perron--Frobenius (still primitive since the
zero-terminating edge $j\mapsto0$ together with the continuation cycle through
$L-1$ remains, using $c_L\ge1$). Dividing by $H_{n+1}-H_n=\Theta(\lambda^n)$
gives the vanishing density. Uniqueness on $\mathcal R_n$ holds because
$\mathcal A^{-}$ is deterministic and reduces the count, so the representation
map is still injective on its image; the Quasi-Power Theorem applied to
$M^{-}(y)$ gives the Gaussian law on the conditioned ensemble.
\end{proof}

\begin{remark}
The trichotomy is governed by comparing the spectral radius $\rho_\bullet$ of the
(possibly modified) transfer matrix with $\lambda$: $\rho_\bullet=\lambda$ gives
unique existence; $\rho_\bullet>\lambda$ gives existence with exponentially many
representations; $\rho_\bullet<\lambda$ gives existence only on an exponentially
sparse set. In all three regimes the number of summands is asymptotically
Gaussian with mean and variance linear in $n$, because in each case the relevant
transfer matrix is primitive with an analytic simple dominant eigenvalue, the
sole hypothesis of the Quasi-Power Theorem; only the constants $(\mu,\sigma^2)$
change.
\end{remark}

\section{Part (b): The limiting gap distribution}\label{sec:b}

We now compute $P(k)=\lim_{n\to\infty}P_n(k)$, the limiting probability that the
gap between the indices of a uniformly chosen consecutive pair of summands equals
$k$, for $N$ uniform in $[H_n,H_{n+1})$. A \emph{gap} of size $k$ between two
consecutive nonzero digits $a_i\neq0,\ a_{i+k}\neq0$ means exactly
$a_{i+1}=\cdots=a_{i+k-1}=0$.

\subsection{The Perron eigenvectors of $M$}

\begin{lemma}[Left and right Perron eigenvectors]\label{lem:eig}
Let $\lambda$ be the dominant root of \eqref{eq:char}. Set
\begin{equation}\label{eq:rj}
r_j=\lambda^{j}-\sum_{i=1}^{j}c_i\,\lambda^{\,j-i}\qquad(0\le j\le L-1),
\quad r_0=1 .
\end{equation}
Then $r=(r_0,\dots,r_{L-1})$ is the left Perron eigenvector of $M$
($M^{\!\top}r=\lambda r$), and $u_j=\lambda^{-j}$ is the right Perron eigenvector
($Mu=\lambda u$). Moreover $r_j>0$ for all $j$.
\end{lemma}

\begin{proof}
$(M^{\!\top}r)_j=\sum_{j'}M_{j'j}r_{j'}=c_{j+1}r_0+[\,j+1<L\,]\,r_{j+1}$.
Imposing $=\lambda r_j$ with $r_0=1$ gives the recursion
$r_{j+1}=\lambda r_j-c_{j+1}$ for $0\le j\le L-2$ and $c_Lr_0=\lambda r_{L-1}$.
Solving the recursion from $r_0=1$ yields \eqref{eq:rj}; the boundary relation
$\lambda r_{L-1}=c_L$ is exactly the characteristic equation \eqref{eq:char}
divided by $\lambda^{0}$, hence holds at $\lambda$. For the right eigenvector,
$(Mu)_j=u_{j-1}=\lambda^{-(j-1)}=\lambda\cdot\lambda^{-j}=\lambda u_j$ for
$1\le j\le L-1$, and $(Mu)_0=\sum_j c_{j+1}\lambda^{-j}=\lambda$ by
\eqref{eq:char}, $=\lambda u_0$. Positivity $r_j>0$ is Perron positivity of the
left eigenvector of the primitive matrix $M$.
\end{proof}

\subsection{The stationary digit chain and the deterministic zero map}

The asymptotic statistics of the digit string of a uniform $N\in[H_n,H_{n+1})$
are governed by the stationary Markov chain on states $\{0,\dots,L-1\}$ with
transition probabilities (the standard Perron / maximal-entropy chain)
\begin{equation}\label{eq:chain}
P(j\to j')=\frac{M_{j'j}\,r_{j'}}{\lambda\,r_j},
\qquad\text{stationary law}\quad
\pi_j=\frac{r_j\,\lambda^{-j}}{S},\quad S=\sum_{i=0}^{L-1}r_i\lambda^{-i}.
\end{equation}
That $P$ is column-stochastic follows from $M^{\!\top}r=\lambda r$, and that
$\pi$ is stationary follows from $Mu=\lambda u$ with $u_j=\lambda^{-j}$
(Lemma~\ref{lem:eig}); both are immediate from the eigenvector relations.
Equation \eqref{eq:density}'s closed form for $S$ follows by substituting
\eqref{eq:rj}:
\[
r_j\lambda^{-j}=1-\sum_{i=1}^{j}c_i\lambda^{-i},\quad
S=\sum_{j=0}^{L-1}\Big(1-\sum_{i=1}^{j}c_i\lambda^{-i}\Big)
 =L-\sum_{i=1}^{L-1}(L-i)c_i\lambda^{-i}.
\]

\begin{lemma}[Convergence to the stationary chain]\label{lem:conv}
For each fixed window pattern $w$ of digits, the empirical frequency of $w$ among
the digits of a uniform $N\in[H_n,H_{n+1})$ converges, as $n\to\infty$, to the
stationary probability of $w$ under \eqref{eq:chain}. Consequently
$P_n(k)\to P(k)$ where $P(k)$ is the stationary probability of the pattern
``nonzero, then $k-1$ zeros, then nonzero'' conditioned on the first digit being
nonzero.
\end{lemma}

\begin{proof}
By Lemma~\ref{lem:transfer}, choosing $N$ uniformly in $[H_n,H_{n+1})$ is the
same as choosing a uniformly random legal length-$n$ path of $\mathcal A$ with
leading digit at index $n$. For a primitive transfer matrix $M$, the empirical
distribution of length-$\ell$ factors of a uniformly random length-$n$ path
converges (as $n\to\infty$, $\ell$ fixed) to the stationary distribution of the
maximal-entropy Markov measure \eqref{eq:chain}; this is the standard
equidistribution theorem for subshifts of finite type / primitive transfer
matrices (Parry measure), whose hypotheses (primitivity, simple dominant
eigenvalue) hold here by Perron--Frobenius. The boundary (first $L$ and last $L$
digits) contributes $O(1/n)$ to any fixed pattern frequency, hence does not
affect the limit. Choosing a uniform consecutive pair of summands and asking for
gap $=k$ is exactly the conditional pattern in the statement, so
$P_n(k)\to P(k)$.
\end{proof}

The crucial simplification is that the zero-transitions of \eqref{eq:chain} form
a \emph{deterministic} map on states.

\begin{lemma}[Deterministic zero map]\label{lem:zero}
Under \eqref{eq:chain}, from state $j$ a digit equal to $0$ forces a unique
destination $f(j)$ with weight $z_j$:
\begin{equation}\label{eq:f}
f(j)=\begin{cases}0,& c_{j+1}\ge1,\\ j+1,& c_{j+1}=0,\end{cases}
\qquad
z_j=\begin{cases}\dfrac{1}{\lambda r_j},& c_{j+1}\ge1,\\[2mm]
\dfrac{r_{j+1}}{\lambda r_j},& c_{j+1}=0.\end{cases}
\end{equation}
In particular $f(0)=0$ with $z_0=1/\lambda$ (as $c_1\ge1$): once the chain
reaches state $0$, every further zero keeps it at $0$ with weight $1/\lambda$.
For any state $s$, following zeros reaches state $0$ in $g(s)\le L-1$ steps (since
$c_L\ge1$ forces a return at state $L-1$), and the accumulated weight telescopes:
\begin{equation}\label{eq:tele}
\prod_{\text{$g(s)$ zero steps from $s$ to $0$}}z_{\bullet}
=\frac{1}{\lambda^{g(s)}\,r_s}.
\end{equation}
\end{lemma}

\begin{proof}
From state $j$ a zero digit is admissible only as either the terminating digit
$a=0$ (present iff $c_{j+1}\ge1$, sending $j\to0$ with probability
$r_0/(\lambda r_j)=1/(\lambda r_j)$) or the continuation digit $a=c_{j+1}=0$
(present iff $c_{j+1}=0$, sending $j\to j+1$ with probability
$r_{j+1}/(\lambda r_j)$). These two cases are mutually exclusive, giving
\eqref{eq:f}. Iterating $f$ from $s$: while $c_{s+1}=c_{s+2}=\cdots=0$ we move
$s\to s+1\to\cdots$ until the first index $m$ with $c_{m+1}\ge1$, where we drop to
$0$; since $c_L\ge1$, $m\le L-1$ and $g(s)=m-s+1\le L-1$. The product of weights
along $s\to s+1\to\cdots\to m\to0$ is
$\frac{r_{s+1}}{\lambda r_s}\cdots\frac{r_m}{\lambda r_{m-1}}\cdot
\frac{1}{\lambda r_m}=\frac{1}{\lambda^{g(s)}r_s}$,
which is \eqref{eq:tele} after cancellation of the telescoping product
$r_{s+1}\cdots r_m$ against $r_s\cdots r_{m-1}$.
\end{proof}

\subsection{The gap distribution}

Let $a_s$ denote the stationary probability of executing a nonzero transition
that \emph{lands} in state $s$, and let $q(t)$ be the probability of emitting a
nonzero digit from state $t$. From \eqref{eq:chain},
\begin{align}
a_{j+1}&=[\,c_{j+1}\ge1,\ j+1<L\,]\;\frac{r_{j+1}}{\lambda r_j}\,\pi_j
&&\text{(nonzero continuation)},\label{eq:a1}\\
a_0&=\sum_{j=0}^{L-1}(c_{j+1}-1)^{+}\,\frac{1}{\lambda r_j}\,\pi_j
&&\text{(nonzero termination, digits $1..c_{j+1}-1$)},\label{eq:a0}\\
q(t)&=[\,c_{t+1}\ge1,\ t+1<L\,]\frac{r_{t+1}}{\lambda r_t}
+(c_{t+1}-1)^{+}\frac{1}{\lambda r_t}.\label{eq:q}
\end{align}
Write $B=\sum_{s}a_s$ for the stationary density of nonzero digits (the mean
number of summands per index; this is exactly $\mu=\beta$ of
Theorem~\ref{thm:gauss}, see Lemma~\ref{lem:base}).

\begin{theorem}[Limiting gap distribution]\label{thm:gap}
For every $k\ge1$,
\begin{equation}\label{eq:Pk}
\boxed{\;
P(k)=\frac{1}{B}\sum_{s=0}^{L-1} a_s\;\Big(\prod_{\text{$k-1$ zero steps from $s$}}z_\bullet\Big)\;q\big(f^{(k-1)}(s)\big)\; }
\end{equation}
with $a_s,q,B,z,f$ given by \eqref{eq:a1}--\eqref{eq:q} and
\eqref{eq:f}. In closed form, $q(0)=(\lambda-1)/\lambda$ universally, and for
all $k>L$ (so that every $s$ has already reached state $0$ after $g(s)\le L-1$
zero steps),
\begin{equation}\label{eq:tail}
P(k)=A\,\lambda^{-k},\qquad
A=\frac{(\lambda-1)}{B}\sum_{s=0}^{L-1}\frac{a_s}{r_s}.
\end{equation}
For $2\le k\le L$ the value is given exactly by \eqref{eq:Pk}; the dependence on
which $c_i$ vanish enters through the deterministic map $f$ in
Lemma~\ref{lem:zero} (a zero at index where $c_{j+1}=0$ moves up the states
rather than resetting to $0$, changing both $f^{(k-1)}(s)$ and the product of
$z$'s).
\end{theorem}

\begin{proof}
By Lemma~\ref{lem:conv}, $P(k)$ is the stationary conditional probability of the
pattern ``nonzero, $(k-1)$ zeros, nonzero''. Condition on the first nonzero:
the state immediately after it is distributed as $a_s/B$. The next $k-1$ digits
are forced to be zero; by Lemma~\ref{lem:zero} this is a deterministic path of
weight $\prod_{i=0}^{k-2}z_{f^{(i)}(s)}$ ending in state $f^{(k-1)}(s)$; the
$k$-th digit is nonzero with probability $q(f^{(k-1)}(s))$. Summing over $s$ and
dividing by the conditioning mass $B$ gives \eqref{eq:Pk}.

For the tail $k>L$: by Lemma~\ref{lem:zero}, after $g(s)\le L-1\le k-1$ zero
steps the chain is at state $0$ and stays there; the accumulated weight is
$\frac{1}{\lambda^{g(s)}r_s}$ (by \eqref{eq:tele}) times $z_0^{\,(k-1)-g(s)}
=\lambda^{-((k-1)-g(s))}$ for the remaining zeros, i.e.
$\frac{1}{\lambda^{g(s)}r_s}\cdot\lambda^{-(k-1-g(s))}
=\frac{1}{r_s\,\lambda^{\,k-1}}$, independent of $g(s)$. The final nonzero is
from state $0$, with probability $q(0)$. Computing
$q(0)=\frac{r_1}{\lambda}+\frac{c_1-1}{\lambda}
=\frac{(\lambda-c_1)+(c_1-1)}{\lambda}=\frac{\lambda-1}{\lambda}$
(using $r_0=1$, $r_1=\lambda-c_1$). Hence
\[
P(k)=\frac{1}{B}\sum_s a_s\,\frac{1}{r_s\lambda^{\,k-1}}\,q(0)
   =\frac{q(0)\,\lambda}{B}\Big(\sum_s\frac{a_s}{r_s}\Big)\lambda^{-k}
   =A\,\lambda^{-k},
\]
which is \eqref{eq:tail}. Validity for all $k>L$ uses $g(s)\le L-1\le k-1$,
established in Lemma~\ref{lem:zero}.
\end{proof}

\begin{lemma}[Density of summands]\label{lem:base}
$B=\beta=1-\frac1S\sum_{j=0}^{L-1}w_j\lambda^{-(j+1)}$ with $w_j$ as in
\eqref{eq:density}; equivalently $B$ is one minus the stationary probability of a
zero digit.
\end{lemma}

\begin{proof}
A step emits a zero with probability $\pi_jz_j$ from state $j$; by
Lemma~\ref{lem:zero}, $\pi_jz_j=\frac{r_j\lambda^{-j}}{S}\cdot
\frac{w_j}{\lambda r_j}=\frac{w_j\lambda^{-(j+1)}}{S}$ where $w_j=1$ if
$c_{j+1}\ge1$ and $w_j=r_{j+1}$ if $c_{j+1}=0$. Summing and subtracting from $1$
gives the formula; $B=\sum_s a_s$ is the complementary nonzero density.
\end{proof}

\subsection{Worked closed forms}
\begin{itemize}
\item \textbf{$L=2$, general $(c_1,c_2)$ with $c_2\ge1$.} Then $r_0=1$,
$r_1=\lambda-c_1$, $S=2-c_1/\lambda$, and \eqref{eq:Pk} simplifies (for all
$k\ge2$) to
\begin{equation}\label{eq:L2}
P(k)=\frac{(\lambda-1)(\lambda^2-1)}{\big(2\lambda^2-(c_1+1)\lambda-1\big)}\;\lambda^{-k}.
\end{equation}
For Fibonacci ($c_1=c_2=1$, $\lambda=\varphi=\frac{1+\sqrt5}{2}$): the prefactor
equals $(\varphi-1)(\varphi^2-1)/(2\varphi^2-2\varphi-1)$; using
$\varphi^2=\varphi+1$ this is $1$, so
\begin{equation}\label{eq:fib}
P(k)=\varphi^{-k}\qquad(k\ge2),\qquad \sum_{k\ge2}\varphi^{-k}=1.
\end{equation}
\item \textbf{$L=3$, $(c_1,c_2,c_3)=(3,2,1)$, $\lambda\approx3.6274$.}
For all $k\ge2$,
\[
P(k)=\frac{3\lambda^2-\lambda-2}{2\lambda^2-4\lambda-1}\;\lambda^{-k}.
\]
\item \textbf{A vanishing coefficient: $(c_1,c_2,c_3)=(2,0,1)$,
$\lambda\approx2.2056$.} Here $c_2=0$ forces $f(1)=2$, so the boundary differs
from the tail:
\[
P(1)=\lambda^{-1},\quad P(2)=\lambda^{-2},\quad P(k)=2\,\lambda^{-k}\ (k\ge3).
\]
\item \textbf{$(c_1,c_2)=(1,2)$, $\lambda=2$.} $P(k)=2^{-k}$ for $k\ge1$.
\item \textbf{Gaps forced larger when interior coefficients vanish:}
for $(c_1,c_2,c_3)=(1,0,1)$, $\lambda\approx1.4656$, the vanishing of $c_2$ makes
$f(1)=2$, forcing $P(1)=P(2)=0$; the deterministic-zero computation
\eqref{eq:Pk} gives $A=1$ and
\[
P(1)=P(2)=0,\qquad P(k)=\lambda^{-k}\quad(k\ge3),
\]
which indeed sums to $\sum_{k\ge3}\lambda^{-k}=\lambda^{-2}/(\lambda-1)=1$ by the
characteristic equation $\lambda^3=\lambda^2+1$.
\end{itemize}

\section{Part (c): The largest gap}\label{sec:c}

Let $G_n=G(N)$ be the largest gap in the legal decomposition of a uniform
$N\in[H_n,H_{n+1})$. The number of summands is $K\sim\beta n$
(Theorem~\ref{thm:gauss}), hence the number of gaps is $K-1\sim\beta n$, and by
Theorem~\ref{thm:gap} each gap is, in the $n\to\infty$ limit, distributed with an
exponential tail $P(k)=A\lambda^{-k}$ ($k>L$). The largest of $\approx\beta n$
such gaps therefore concentrates at scale $\log_\lambda n$, and after the natural
discrete centering converges to a Gumbel (doubly-exponential) law.

\begin{theorem}[Largest-gap law]\label{thm:max}
Let $\lambda$ be the Perron root and let $A,B=\beta$ be the constants of
Theorem~\ref{thm:gap} and Lemma~\ref{lem:base}. Define the centering
\begin{equation}\label{eq:center}
b_n=\log_\lambda\!\big(A\,\beta\,n\big)
   =\frac{\log\!\big(A\beta n\big)}{\log\lambda}.
\end{equation}
Then for every real $x$ and every sequence of integers $k_n=\lfloor b_n+x\rfloor$,
\begin{equation}\label{eq:gumbel}
\lim_{n\to\infty}\Prob\!\big(G_n\le b_n+x\big)
=\exp\!\big(-\lambda^{-\{b_n+x\}}\,\lambda^{-\lfloor x\rfloor}\big)
=\exp\!\big(-\lambda^{-(b_n+x-\lfloor b_n+x\rfloor)}\,A\beta n\,\lambda^{-(b_n+x)}\big),
\end{equation}
i.e.\ $G_n-b_n$ converges, along the lattice and modulo the
$O(1)$-periodic fluctuation $\{b_n\}$ inherent to extreme values of integer
geometric variables, to the discrete Gumbel distribution with cumulative
distribution function $x\mapsto \exp(-\lambda^{-x})$ on the real line. In
particular
\begin{equation}\label{eq:meanmax}
\E[G_n]=\log_\lambda(\beta n)+O(1),\qquad \mathrm{Var}(G_n)=O(1),
\end{equation}
and $G_n=(1+o(1))\dfrac{\log n}{\log\lambda}$ almost surely along the ensemble.
\end{theorem}

\begin{proof}
Fix $t\ge1$ and let $E_n(t)=\Prob(G_n<t)$ be the probability that \emph{no} gap
is $\ge t$. By Lemma~\ref{lem:conv} the digit string of a uniform
$N\in[H_n,H_{n+1})$ converges to the stationary chain of \eqref{eq:chain}, whose
gaps between consecutive nonzeros, by Theorem~\ref{thm:gap}, form a stationary
sequence with marginal $P(\cdot)$ and an exponential tail
$\Prob(\text{gap}\ge t)=\sum_{k\ge t}A\lambda^{-k}
=\frac{A}{\lambda-1}\lambda^{-(t-1)}=:c_0\lambda^{-t}$ for $t>L$
(with $c_0=A\lambda/(\lambda-1)$). The number of gaps is
$K-1=\beta n(1+o(1))$ with exponentially small deviations
(Theorem~\ref{thm:gauss}). Although consecutive gaps are not independent, the
stationary gap process is exponentially $\phi$-mixing: the gap chain is a
function of the finite-state primitive chain \eqref{eq:chain}, whose correlations
decay at rate equal to the spectral gap $|\lambda_2/\lambda|<1$ of $M$. For such
mixing stationary sequences the classical extreme-value theory (Leadbetter's
conditions $D(u_n)$ and $D'(u_n)$, both verified here because the exceedance
events $\{\text{gap}\ge u_n\}$ have probability $\to0$ and the mixing is
exponential) gives the Poisson approximation: the number of gaps $\ge u_n$ among
the $\approx\beta n$ gaps converges to a Poisson random variable with mean
$\Lambda_n=\beta n\cdot c_0\lambda^{-u_n}$, hence
\[
\Prob(G_n<u_n)=\exp(-\Lambda_n)+o(1)
=\exp\!\big(-\beta n\,c_0\,\lambda^{-u_n}\big)+o(1).
\]
Choosing $u_n=b_n+x$ with $b_n$ as in \eqref{eq:center} makes
$\beta n\,c_0\,\lambda^{-u_n}
=\frac{A\lambda}{\lambda-1}\lambda^{-x}\cdot\lambda^{-\{b_n\}}$
times a bounded periodic factor, yielding \eqref{eq:gumbel}: the limit is
$\exp(-c\,\lambda^{-x})$ for the appropriate constant $c$, a Gumbel law on the
$\log_\lambda$ scale, with the lattice/periodic correction $\{b_n\}$ standard for
maxima of integer-valued geometric variables. The mean and variance
\eqref{eq:meanmax} follow from the Gumbel limit: the location is
$b_n=\log_\lambda(\beta n)+O(1)$ and the spread is $O(1)$ since the Gumbel scale
parameter is $1/\log\lambda$, a constant. The almost-sure first-order asymptotics
$G_n\sim\log n/\log\lambda$ is the Borel--Cantelli consequence of
$\Prob(|G_n-\log_\lambda n|>\varepsilon\log n)\to0$ exponentially.
\end{proof}

\begin{remark}[Dependence on the $c_i$]
The largest-gap law depends on $(c_1,\dots,c_L)$ only through three explicit
constants: the Perron root $\lambda$ (the decay rate of the gap tail, controlling
the spread $1/\log\lambda$ and the centering scale), the tail amplitude $A$ of
\eqref{eq:tail}, and the summand density $\beta=B$ of Lemma~\ref{lem:base}
(controlling the number of gaps). When some interior $c_i$ vanish, the smallest
possible gap increases (e.g.\ all gaps are $\ge2$ for Fibonacci, and $\ge3$ when
$c_2=0$), but the \emph{largest} gap is governed solely by the common tail
$A\lambda^{-k}$, so \eqref{eq:gumbel}--\eqref{eq:meanmax} hold verbatim with the
$(\lambda,A,\beta)$ computed for the given coefficients.
\end{remark}

\section*{Summary of answers}
\begin{itemize}
\item[(a)] Existence/uniqueness, mean summands, and Gaussian fluctuations are
governed by the spectral radius $\rho_\bullet$ of the (possibly modified)
transfer matrix relative to $\lambda$. Unique existence $\Leftrightarrow
\rho_\bullet=\lambda$; relaxing legality to $\rho_\bullet>\lambda$ produces
$\Theta((\rho_\bullet/\lambda)^n)$ legal decompositions on average
(Prop.~\ref{prop:multi}); tightening to $\rho_\bullet<\lambda$ leaves only an
exponentially sparse representable set (Prop.~\ref{prop:none}). In all regimes
the number of summands is asymptotically normal with mean and variance linear in
$n$ (Quasi-Power Theorem), only the constants changing.
\item[(b)] The limiting gap distribution is \eqref{eq:Pk}: with the explicit
Perron data $r_j=\lambda^j-\sum_{i\le j}c_i\lambda^{j-i}$, $\pi_j\propto
r_j\lambda^{-j}$, the deterministic zero map \eqref{eq:f}, and
$a_s,q,B$ of \eqref{eq:a1}--\eqref{eq:q}. It has the universal exponential tail
$P(k)=A\lambda^{-k}$ for $k>L$ with $A=\frac{\lambda-1}{B}\sum_s a_s/r_s$, while
for $2\le k\le L$ it is given by \eqref{eq:Pk} with the boundary depending on
which $c_i$ vanish through $f$. For Fibonacci, $P(k)=\varphi^{-k}$.
\item[(c)] The largest gap satisfies $G_n=\log_\lambda(\beta n)+O(1)$ and, after
centering by $b_n=\log_\lambda(A\beta n)$, converges to a Gumbel law on the
$\log_\lambda$ scale, \eqref{eq:gumbel}, with all dependence on
$(c_1,\dots,c_L)$ through $\lambda$, $A$, and $\beta=B$.
\end{itemize}

\end{document}
